\section{Introduction}
\label{sec:intro}

% [TODO prose] Expand each paragraph; this is a structured draft with the key claims.

Reservoirs regulate water supply, irrigation, hydropower and flood risk, and their
management depends on knowing how much water they store as a function of level,
i.e.\ the area--elevation--volume (AEV) relationship. For many dams the operational
AEV curve dates to construction, decades ago, and no longer reflects the basin's
present geometry after prolonged sedimentation \citep{BrunoEtAl2003}. Repeating the
bathymetric surveys that produced those curves is costly and infrequent, so most
reservoirs are monitored, if at all, with a stage record and an obsolete curve.

Satellite remote sensing offers a route to updating this information without field
campaigns. Two complementary signals are available. Satellite radar altimetry, as
compiled in databases such as DAHITI \citep{Schwatke2015}, and now the wide-swath
mission SWOT \citep{Biancamaria2016}, measures water-surface elevation; Synthetic
Aperture Radar (SAR) such as Sentinel-1 \citep{Torres2012} measures water-surface
\emph{area} directly, independently of any altimetry track and with a reliability
that a companion study shows is predictable a priori from shoreline geometry
\citep{PaperOne}. \citet{Schwatke2020} combined optical/SAR water masks with
altimetry to reconstruct hypsometry for large lakes. Here we adapt that idea to
the reservoir-monitoring problem and, crucially, to the many reservoirs that have
\emph{neither} an in-situ gauge \emph{nor} a nadir altimetry track over them.

Our premise is that the shoreline exposed as a reservoir draws down is a sequence
of contour lines at known water levels: stacking Sentinel-1 waterlines ordered by
level reconstructs a digital elevation model (DEM) of the emerged basin, from which
an updated AEV curve and storage follow by integration. The water level needed to
label each waterline is the linchpin. Where a gauge exists it is trivial; where it
does not, the level must itself come from remote sensing. We therefore test whether
the in-situ gauge can be replaced, first by nadir altimetry (DAHITI) and then by
SWOT swath altimetry, which observes reservoirs regardless of track and so answers
the ungauged, track-missed case.

The contributions of this paper are: (i) a fully remote-sensing pipeline that
reconstructs updated reservoir bathymetry and storage from Sentinel-1 waterlines
and a remote water level; (ii) a validation against a December~2025 echo-sounder
survey and recent official bathymetric curves on nine Sicilian reservoirs;
(iii) a water-level-substitution experiment showing that SWOT altimetry yields
storage equivalent to an in-situ gauge and to field truth; (iv) an uncertainty and
observability analysis linking reconstruction reliability to shoreline
area-to-perimeter ratio and to drawdown amplitude; and (v) an open, interactive
tool that runs the pipeline.
